% ID: SEL-D0011-BD-E08
% ESTADO: cotejada. Segunda pasada visual completa y borrador compilado revisados.
% Fuente: Selectividad/matematicas_aplicadas/EXTRA_SUPLENTE_Modelo-1-Enunciados (Mat-Ap-2021).pdf
% SHA-256: 4a6c7d34ff35c3e7678c52970b4d4558cee2c35224b46c66035cb0bc8d5a7ec8
% Página: 2; bloque D; ejercicio original 8.

El tiempo, en horas, que los alumnos de un instituto dedican a estudiar
para los exámenes finales, se distribuye siguiendo una ley Normal de media
desconocida y varianza 81. Se toma una muestra aleatoria de 16 alumnos de dicho
instituto, obteniéndose los siguientes tiempos:
\[
30\quad42\quad38\quad45\quad52\quad60\quad21\quad26\quad
33\quad44\quad28\quad49\quad32\quad51\quad49\quad40
\]
\begin{enumerate}[label=\alph*)]
  \item \textbf{(1.5 puntos)} Obtenga un intervalo, con un 95\% de confianza,
  para estimar el tiempo medio de estudio de los alumnos de ese instituto.
  \item \textbf{(1 punto)} Calcule el mínimo tamaño de la muestra que se ha
  de tomar, para estimar el tiempo medio de estudio de esos alumnos con un
  error inferior a 2 horas y un nivel de confianza del 98\%.
\end{enumerate}
